\begin{frame}
	\frametitle{RNA vs. RNP: o par que se confunde}
	\small
	\begin{columns}
		\column{.48\linewidth}
		\begin{beamercolorbox}[sep=5pt,rounded=true,shadow=true]{title}
			Rede Neural Artificial (RNA)
		\end{beamercolorbox}
		\begin{itemize}
			\item Ativação: valor \textbf{contínuo} por neurônio, por camada.
			\item Processamento \textbf{síncrono}: uma passada completa por amostra.
			\item Sem dimensão temporal intrínseca (a menos que seja recorrente).
			\item Custo computacional fixo, independente da atividade.
		\end{itemize}
		\column{.48\linewidth}
		\begin{beamercolorbox}[sep=5pt,rounded=true,shadow=true]{title}
			Rede Neural de Pulso (RNP)
		\end{beamercolorbox}
		\begin{itemize}
			\item Ativação: \textbf{pulso} (0 ou 1), esparso no tempo.
			\item Processamento \textbf{orientado a evento}: só computa quando há pulso.
			\item Dinâmica temporal \textbf{intrínseca}: potencial de membrana acumula ao longo do tempo.
			\item Informação codificada em \textbf{taxa} e/ou \textbf{tempo} de disparo \cite{kasabov2019time}.
		\end{itemize}
	\end{columns}
	\vspace{2pt}
	\begin{alertblock}{Consequência prática}
		Uma RNP processa naturalmente \textbf{fluxos} de dados (áudio, EEG, eventos de câmera) --- mas processar um único dado estático exige ``esticá-lo'' ao longo de vários passos de tempo, um custo que a RNA não tem.
	\end{alertblock}
\end{frame}
